\subsection{KT-TRACE-AVOID: the exact single-element obstruction}
\label{sec:kt-trace-avoid}

Let $\Gamma$ be a profinite group and let
$\overline\rho_1,\ldots,\overline\rho_m,\overline T$
be continuous finite-dimensional representations over
$\overline{\mathbf F}_\ell$, with finite image. They
factor through their common finite image $G$. In the
application $\overline T=\overline\Psi\oplus1\oplus
\overline\Psi^{-1}$ and all dimensions are three.
Write $K_0(G)$ for the Grothendieck ring of finite-dimensional
$\overline{\mathbf F}_\ell[G]$-modules, using tensor product
as multiplication, and form
\[
        V=\prod_{i=1}^m([\overline\rho_i]-[\overline T])
                          \quad\hbox{in }K_0(G).
                                                               \tag{TA.1}
\]

\begin{proposition}[KT-TRACE-AVOID]
There is an element $\sigma\in\Gamma$ such that
\[
 \operatorname{tr}\overline\rho_i(\sigma)
       \ne\operatorname{tr}\overline T(\sigma)
                     \qquad\hbox{for every }i
                                                               \tag{TA.2}
\]
if and only if $V\notin\ell K_0(G)$. Thus the obstruction
is an exact virtual tensor-product relation modulo $\ell$;
it is not in general equivalent to one factor vanishing.
\end{proposition}

\begin{proof}
Let $S_1,\ldots,S_r$ be the simple modules over the finite
group algebra $k[G]$, where $k=\overline{\mathbf F}_\ell$.
Their trace functions are linearly independent over $k$.
Indeed a relation on all group elements extends linearly
to the whole group algebra. Its semisimple quotient is
$\prod_j\operatorname{End}_k(S_j)$, and a lift of a single
diagonal matrix unit in the $j$th factor forces the $j$th
coefficient of the relation to be zero. This argument
does not require $\ell\nmid|G|$.

Trace is additive in short exact sequences and multiplicative
under tensor product. Therefore the trace of $V$ at $g$ is
\[
 \operatorname{tr}V(g)=
 \prod_i\bigl(\operatorname{tr}\overline\rho_i(g)
                   -\operatorname{tr}\overline T(g)\bigr).
                                                               \tag{TA.3}
\]
Since $K_0(G)$ is free on the classes $[S_j]$, the trace
function of $V$ vanishes identically exactly when every
integer coefficient of $V$ in this basis is divisible by
$\ell$, by the independence just proved. Equation (TA.3)
is nonzero at some $g$ exactly when every factor there is
nonzero. Lift $g$ to $\Gamma$ to obtain (TA.2).
\end{proof}

\paragraph{Individual traces versus simultaneous avoidance.}
For two semisimple representations of the same dimension
$n<\ell$, equality of traces on every group element implies
isomorphism. The preceding independence proves this directly:
each difference of simple multiplicities lies in $[-n,n]$,
so divisibility by $\ell$ forces it to be zero. Thus for
$n=3$, $\ell>3$, a single factor in (TA.1) vanishes modulo
$\ell$ exactly when its residual semisimplification is
$\overline T$. Nevertheless a product of nonzero factors
can vanish, as the following explicit example shows.

\begin{example}[KT-TRACE-COVER]
For every prime $\ell>3$, take
$G=C_3\times C_2\times C_2$ and choose a faithful character
$\Psi:C_3\to\overline{\mathbf F}_\ell^\times$.
Let $\chi_1,\chi_2,\chi_3$ be the three nontrivial
characters of $C_2\times C_2$, extended trivially on $C_3$.
Set
\[
 T=\Psi\oplus1\oplus\Psi^{-1},\qquad
 R_i=(\Psi\chi_i)\oplus1\oplus(\Psi^{-1}\chi_i).
                                                               \tag{TA.4}
\]
The three $R_i$ are pairwise nonisomorphic semisimple
three-dimensional representations; none is isomorphic to
$T$, and $\Psi^2\ne1$. Each has determinant one and is
self-dual. For $g=(h,v)$,
\[
 \operatorname{tr}R_i(g)-\operatorname{tr}T(g)
       =(\chi_i(v)-1)(\Psi(h)+\Psi(h)^{-1}).
                                                               \tag{TA.5}
\]
Every $v\in C_2\times C_2$ lies in the kernel of at least
one nontrivial $\chi_i$. Hence every $g$ makes at least
one difference in (TA.5) zero; no single element separates
all three systems. These statements hold for arbitrarily
large $\ell$ while the number and dimensions of the
systems stay fixed. This is an abstract residual example,
not an assertion that these representations are members
of the automorphic admissible class (GR.1).
\end{example}

\begin{proof}
Choose independent quadratic characters $\chi_1,\chi_2$
and put $\chi_3=\chi_1\chi_2$. The three nonzero elements
of $C_2\times C_2$ each have one of these characters equal
to $1$ and the other two equal to $-1$; at the identity
all three are $1$. Their kernels therefore cover the
group. The characters occurring in (TA.4) are distinct
because $\Psi$ has order three and the $\chi_i$ have
order two. These observations prove all the assertions,
including the displayed trace formula.
\end{proof}

This example also illustrates the stronger multi-operator
criterion rechecked alongside K17-S17-JOINT. At $h=1$, choose nonzero
$v_1,v_2$ with respectively $\chi_1(v_1)=1$ and
$\chi_2(v_2)=1$. On the three-system eigenvalue space,
the two trace-difference operators are
\[
 A_1=\operatorname{diag}(0,-4,-4),\qquad
 A_2=\operatorname{diag}(-4,0,-4).
                                                               \tag{TA.6}
\]
Each determinant is zero, but $[A_1\mid A_2]$ has full
row rank three for every odd $\ell$. Thus a multi-operator
certificate can succeed even when every individual group
element fails. A linear combination of group elements is
not itself a Galois element, so this does not change the
literal single-element assertion of \texttt{OneCharU\_s}.

